\itemhead{复习题二 · 第9题（教材p.62）}{XJ-B1-C02-RV-P062-Y0642}
\label{XJ-B1-C02-RV-P062-Y0642}
{\footnotesize 湘教版教材p.62 / PDF第69页 / 复习题二 / 第9题}
\field{原题}{若 $x>0,y>0$，且 $xy=x+y+3$，求 $x+y$ 的取值范围。}
\field{参考答案}{设 $s=x+y>0$，$s+3=xy\le s^2/4$，故 $s\ge6$；反之任意 $s\ge6$ 可由正根和为 $s$、积为 $s+3$ 实现，所以取值范围为 $[6,+\infty)$。}
{\small \field{关联知识}{I03 基本不等式、等号与最值建模（教材帮p.52、53、54 / PDF61、62、63）}
\field{方法与易错点}{列正数条件与固定量，构造基本不等式，解等号条件并代回实际限制。 易错：把一个下界误当最小值；多个等号条件不能同时实现；漏单位和整数限制。}
\field{答案依据}{答案册相关题组：PDF第15、16页（册内页14、15）；按当前题库原题逐项重算，不按旧版题号机械对套。}}
